\chapter{บทนำ}
\label{chapter:introduction}

\section{ที่มาและความสำคัญ}

    ในปัจจุบัน เทคโนโลยี AI ได้เข้ามามีบทบาทสำคัญในวงการแพทย์ โดยจากโครงงานการพัฒนาระบบวิเคราะห์รอยโรคทางช่องปากก่อนหน้านี้ ได้ประสบความสำเร็จในการสร้างแบบจำลอง AI ที่สามารถวิเคราะห์โรคได้อย่างแม่นยำ ซึ่งถือเป็นรากฐานที่ดีเยี่ยมในการช่วยทันตแพทย์คัดกรองผู้ป่วย

    อย่างไรก็ตาม ระบบดังกล่าวยังอยู่ในระดับ "งานวิจัยต้นแบบ" เมื่อพิจารณาในมุมมองของวิศวกรรมซอฟต์แวร์ พบว่าระบบเดิมยังมีข้อจำกัดในการนำไปให้บริการจริง เช่น โครงสร้างที่กินทรัพยากรคลาวด์สูงทำให้มีต้นทุนแพง ขาดระบบการจัดการวงจรชีวิตซอฟต์แวร์ ที่ทำให้การอัปเดตระบบทำได้ยาก ตลอดจนยังขาดระบบยืนยันตัวตนและการจัดการข้อมูลที่ปลอดภัยตามมาตรฐาน PDPA

    ด้วยเหตุนี้ โครงงานนี้จึงเกิดขึ้นเพื่อต่อยอดและปิดช่องโหว่ดังกล่าว โดยมุ่งเน้นการปรับปรุงสถาปัตยกรรมระบบใหม่ทั้งหมด นำหลักการ CI/CD มาใช้ ปรับแต่งโมเดล AI ให้ทำงานเร็วและประหยัดต้นทุนขึ้น รวมถึงสร้าง Dashboard เพื่อให้แพทย์สามารถตรวจสอบผลและส่ง Feedback กลับสู่ระบบได้ ซึ่งจะช่วยเปลี่ยนโปรเจกต์นี้ให้กลายเป็นซอฟต์แวร์ทางการแพทย์ที่เสถียร ปลอดภัย และพร้อมใช้งานจริงอย่างยั่งยืน

\section{วัตถุประสงค์}

    1.2.1 เพื่อพัฒนาระบบผู้เชี่ยวชาญวิเคราะห์รอยโรคทางช่องปากให้เป็นแพลตฟอร์มที่เสถียรและพร้อมใช้งานจริง โดยประยุกต์ใช้วงจรชีวิตการพัฒนาซอฟต์แวร์ และการนำส่งระบบอัตโนมัติ (CI/CD)
    
    1.2.2 เพื่อสร้างสภาพแวดล้อมที่ปลอดภัยสำหรับการประมวลผลข้อมูลรอยโรคผู้ป่วย โดยพัฒนาระบบรักษาความปลอดภัยและการยืนยันตัวตนให้สอดคล้องกับมาตรฐานการคุ้มครองข้อมูลส่วนบุคคล (PDPA)
    
    1.2.3 เพื่อเพิ่มประสิทธิภาพและลดต้นทุน การทำงานของแบบจำลองปัญญาประดิษฐ์วิเคราะห์รอยโรคเดิม ผ่านเทคนิค Model Optimization และการปรับสถาปัตยกรรมคลาวด์ให้รองรับผู้ใช้งานจำนวนมากแบบคู่ขนาน
    
    1.2.4 เพื่อสร้างเว็บแอปพลิเคชันสนับสนุนการตัดสินใจสำหรับทันตแพทย์ ที่อำนวยความสะดวกในการตรวจสอบรอยโรค พร้อมระบบให้ข้อมูลย้อนกลับ (Feedback Loop) สำหรับรวบรวมข้อมูลไปพัฒนาความแม่นยำของโมเดลในอนาคต

\section{วิธีการดำเนินงาน}

    1.3.1 ศึกษาและวิเคราะห์โครงสร้างระบบเดิม และ รวบรวมข้อมูลภาพถ่ายช่องปาก ข้อมูลจากการซักประวัติและผลวินิจฉัยจากผู้เชี่ยวชาญ

    1.3.2 ออกแบบสถาปัตยกรรมระบบใหม่ (Microservices) เพื่อรองรับการทำงานที่เสถียรยิ่งขึ้น
    
    1.3.3 นำแบบจำลองเดิมมาทำการปรับแต่ง (Model Optimization) ให้อยู่ในรูปแบบที่เหมาะสม เพื่อลดขนาดและเพิ่มความเร็วในการอนุมานผล

    1.3.4 พัฒนา Web Application และ Doctor Dashboard สำหรับทันตแพทย์โดยให้ทันตแพทย์สามารถกดยืนยันหรือแก้ไขผลการทำนายของ AI ได้เพื่อนำข้อมูลไปพัฒนาความแม่นยำต่อไปในอนาคต

    1.3.5 จัดทำ Automated Testing สำหรับตรวจสอบ API และทดสอบความปลอดภัยของระบบ (Security/Vulnerability Scanning) รวมถึงการทำ Load Testing

    1.3.6 ออกแบบและติดตั้งระบบ CI/CD Pipeline ควบคู่กับการติดตั้งระบบ Centralized Logging เพื่อตรวจจับข้อผิดพลาด (Error Tracking) ในการนำส่งและดูแลรักษาระบบ

    1.3.7 สรุปผลการดำเนินงานของระบบ

\section{ขอบเขตโครงงาน}

    1.4.1  พัฒนาส่วนติดต่อผู้ใช้งานสำหรับทันตแพทย์ เพื่อใช้ในการกรอกข้อมูลซักประวัติ อัปโหลดรูปภาพรอยโรค และแสดงผลการวิเคราะห์ความเสี่ยงจากแบบจำลอง AI 
    
    1.4.2 ออกแบบและพัฒนาระบบหลังบ้าน ในรูปแบบ Microservices บนสภาพแวดล้อม Cloud เพื่อให้รองรับการใช้งานพร้อมกันของคลินิกหรือทันตแพทย์หลายราย

    1.4.3 พัฒนาระบบยืนยันตัวตน ระบบจัดการสิทธิ์ผู้ใช้งาน และการเข้ารหัสข้อมูล เพื่อให้การจัดเก็บข้อมูลผู้ป่วยสอดคล้องกับพ.ร.บ. คุ้มครองข้อมูลส่วนบุคคล (PDPA)

    1.4.4 พัฒนาระบบที่อนุญาตให้ทันตแพทย์สามารถกดยืนยัน หรือแก้ไขผลการวินิจฉัยของ AI ได้ เพื่อรวบรวมข้อมูล Ground Truth กลับไปใช้พัฒนาความแม่นยำของโมเดลในอนาคต

    1.4.5 ติดตั้งระบบ Continuous Integration / Continuous Deployment ควบคู่กับการทำ Automated Testing เพื่อให้การอัปเดตระบบเป็นไปอย่างรวดเร็วและปลอดภัย

    1.4.6  ติดตั้งระบบ Centralized Logging เพื่อตรวจสอบการทำงานของระบบและแจ้งเตือนข้อผิดพลาด แบบเรียลไทม์

\section{ประโยชน์ที่คาดว่าจะได้รับ}

    1.5.1 สามารถยกระดับจากงานวิจัยต้นแบบ ให้กลายเป็นแพลตฟอร์มที่เสถียร รองรับผู้ใช้งานจำนวนมาก และสามารถนำไปติดตั้งใช้งานจริงในคลินิกหรือโรงพยาบาลได้

    1.5.2 สถาปัตยกรรมระบบและโมเดลที่ได้รับการ Optimize จะช่วยลดค่าใช้จ่ายในการเช่าทรัพยากร Cloud Server ลงได้อย่างมีนัยสำคัญ

    1.5.3 ผู้ใช้งานและผู้ป่วยสามารถมั่นใจได้ว่าข้อมูลสุขภาพส่วนบุคคลจะถูกจัดเก็บและประมวลผลอย่างปลอดภัยตามมาตรฐาน PDPA

    1.5.4 ช่วยเพิ่มประสิทธิภาพในการบริหารจัดการเอกสารและการสร้างผลลัพธ์ของงานเมื่อผ่านการอนุมัติครบตามกระบวนการ

    1.5.5 การมีระบบ Feedback Loop จะช่วยให้ทีมผู้พัฒนาโมเดลมีชุดข้อมูลที่ผ่านการคัดกรองจากผู้เชี่ยวชาญ (แพทย์) เพื่อนำไป Train โมเดลให้แม่นยำยิ่งขึ้นในเวอร์ชันถัดไป

    1.5.6 ช่วยลดระยะเวลาในการวิเคราะห์เบื้องต้น ทำให้ทันตแพทย์สามารถคัดกรองรอยโรคที่อาจนำไปสู่โรคมะเร็งช่องปากได้รวดเร็วและเป็นระบบมากขึ้น

    1.5.7 ได้องค์ความรู้และประสบการณ์ในการออกแบบและพัฒนาระบบที่ผสานการจัดการข้อมูล ระบบ Workflow ระบบอนุมัติ และ AI เข้าด้วยกัน

    1.5.8 ได้ระบบต้นแบบที่สามารถต่อยอดหรือปรับใช้กับกระบวนการงานอื่นภายในสถาบันในอนาคตได้
